% !PRINTED none
% Part-title / section-opening page: no folio printed. Centred divider "Twelfth
% Chapter. / Orders in the quadratic field" with a footnote marker, a rule,
% then the "§ 96. Discriminants of the orders." heading and the start of running
% text (Definition 1, remarks 1.-2., footnote). Sequence implies printed page 351.

\begin{center}
\gesperrt{Twelfth Chapter.}

Orders in the quadratic field\textsuperscript{1}).
\end{center}

\rule{3cm}{0.4pt}

\begin{center}
\textbf{\S\ 96. \gesperrt{Discriminants of the orders.}}
\end{center}

1. \gesperrt{Definition. If $Q$ is a natural number, then the totality of the integers of the quadratic field $\Omega$ that are congruent to a rational number with respect to the modulus $Q$ is called an order, and $Q$ is called the conductor of this order.} The order with the conductor $Q$ is denoted by $[Q]$. From this definition it follows:

1. \gesperrt{All rational integral numbers belong to every order;}

2. \gesperrt{Sum, difference and product of two numbers of the order belong to the same order;}

i.\,e. within the order one can carry out addition, subtraction and multiplication unconditionally, but not division in general.

Let the totality of all \gesperrt{integers} of the field $\Omega$ likewise form an order (of conductor 1), the \gesperrt{maximal order}, which would accordingly be denoted by $[1]$.

\rule{4cm}{0.4pt}

\textsuperscript{1}) \gesperrt{Dedekind} introduced the expression order in the general theory of algebraic numbers, because in the special case of the quadratic field they correspond to the \gesperrt{Gauss} orders of quadratic forms. \gesperrt{Hilbert} uses for it the expression \glqq Ring\grqq{} or \glqq Zahlring\grqq{}. Cf. \gesperrt{Dirichlet-Dedekind}, Vorlesungen über Zahlentheorie (\S\ 172 of the third, \S\ 170 of the fourth edition). \gesperrt{Dedekind}: Über die Anzahl der Idealklassen in den verschiedenen Ordnungen eines endlichen Körpers, Festschrift zur Säkularfeier von \gesperrt{Gauss'} Geburtstag, Braunschweig 1877. \gesperrt{Hilbert}: Die Theorie der algebraischen Zahlen, Bericht der Deutschen Mathematischen Vereinigung von 1894/95, Kap. IX.
